\lysh{15688}{4}
\begin{alist}
\item Η συνάρτηση ορίζεται μόνο όταν $e^x - 1 > 0$. Είναι:
\[e^x - 1 > 0 \Leftrightarrow e^x > 1 \Leftrightarrow e^x > e^0 \Leftrightarrow x > 0.\]
Άρα $A = (0, +\infty)$. \\
Η τετμημένη του κοινού σημείου της γραφικής παράστασης $C_f$ της $f$ με τον $x'x$ είναι η λύση της εξίσωσης $f(x) = 0$. Είναι:
\begin{align*}
f(x) = 0 &\Leftrightarrow \ln(e^x - 1) = 0 \Leftrightarrow \ln(e^x - 1) = \ln 1 \Leftrightarrow e^x - 1 = 1 \\
&\Leftrightarrow e^x = 2 \Leftrightarrow x = \ln 2
\end{align*}
Επομένως η $C_f$ τέμνει τον άξονα $x'x$ στο σημείο $(\ln 2, 0)$.
\item Με $x > 0$ έχουμε:
\begin{align*}
f(x) = x - 1 &\Leftrightarrow \ln(e^x - 1) = x - 1 \Leftrightarrow \ln(e^x - 1) = \ln e^{x - 1} \Leftrightarrow e^x - 1 = e^{x - 1} \\
&\Leftrightarrow e^{x + 1} - e = e^x \Leftrightarrow (e - 1)e^x = e \Leftrightarrow e^x = \frac{e}{e - 1} \Leftrightarrow x = \ln\frac{e}{e - 1} \Leftrightarrow x = 1 - \ln(e - 1)
\end{align*}
που περιέχεται στο $A = (0, +\infty)$, αφού
\[e > e - 1 \Rightarrow \ln e > \ln(e - 1) \Rightarrow 1 - \ln(e - 1) > 0\]
\item Έστω $\alpha > 0$. Αν υποθέσουμε ότι η γραφική παράσταση της $f$ έχει κοινά σημεία με την ευθεία $y = x + \alpha$, τότε η εξίσωση $f(x) = x + \alpha$ έχει λύση στο $A$. Είναι:
\[f(x) = x + \alpha \Leftrightarrow \ln(e^x - 1) = \ln e^{x + \alpha} \Leftrightarrow e^x - 1 = e^{x + \alpha} \Leftrightarrow e^{x + \alpha} - e^x = -1\]
που είναι άτοπο, αφού με $\alpha > 0$ ισχύει
\[x + \alpha > x \Rightarrow e^{x + \alpha} > e^x \Rightarrow e^{x + \alpha} - e^x > 0.\]
Επομένως, αν $\alpha > 0$ τότε η γραφική παράσταση της $f$ δεν έχει κοινά σημεία με την ευθεία $y = x + \alpha$.
\end{alist}

